\documentclass[9pt,twocolumn]{extarticle}
\usepackage[margin=1.6cm,columnsep=0.6cm]{geometry}
\usepackage{graphicx}
\usepackage{amsmath,amssymb}
\usepackage{booktabs}
\usepackage{multirow}
\usepackage{hyperref}
\usepackage{xcolor}
\usepackage{caption}
\usepackage{times}
\hypersetup{colorlinks=true,linkcolor=blue,citecolor=blue,urlcolor=blue}
\captionsetup{font=footnotesize,labelfont=bf}

\title{\vspace{-1.2cm}\textbf{¿QEM, QEC o FTQC? Un marco de decisión empírico basado en\\
técnicas de reducción de circuitos y su validación contra el estado del arte}}
\author{\footnotesize Sesión de investigación en cómputo cuántico \\[-2pt]
\footnotesize \texttt{alberto.maldo1312@gmail.com}}
\date{\footnotesize 26 de septiembre de 2026}

\begin{document}
\maketitle

\begin{abstract}
\noindent
Decidir si un circuito cuántico dado debe ejecutarse bajo mitigación de
errores cuánticos (QEM), corrección de errores (QEC) o tolerancia a fallos
completa (FTQC) depende críticamente de qué técnica de reducción/optimización
de circuito se aplique antes de esa decisión --- una dependencia que rara vez
se hace explícita. Este trabajo construye y valida empíricamente un
\emph{pipeline} de decisión (\texttt{pipeline\_v2}/\texttt{decide\_regime})
que normaliza todas las técnicas comparadas al mismo \emph{basis gate set}
real de un backend antes de medir T-count, conteo de compuertas de 2 qubits y
profundidad, evitando así una clase de error metodológico (T-count
artificialmente nulo) que documentamos y corregimos. Sobre 4 familias de
circuito y 3 técnicas de preparación (sin reducir, \texttt{transpile}
optimization\_level=3, cálculo-ZX \texttt{full\_reduce}) encontramos un caso
límite genuino en el que el cálculo-ZX hace que FTQC gane en tiempo esperado
(hasta 212$\times$), aunque la restricción de presupuesto de qubits físicos
sigue bloqueando FTQC en términos absolutos. Extendemos el análisis a seis
técnicas adicionales de reducción basadas en mediciones intermedias y
estructura de circuito (des-cómputo por medición de puertas AND, propagación
de constantes cuántica acotada validada contra BQCP, corte de circuitos,
RQAOA, y enrutamiento por \texttt{SwapStrategy}), validamos tres métricas
propias contra números publicados (dentro de 10--25 puntos porcentuales en
todos los casos comparables), y reportamos honestamente dos resultados
negativos: que encadenar cálculo-ZX con teletransportación de puertas
\emph{empeora} la conectividad del circuito en 3 de 4 casos, y que la
intuición de "qubits de control con mayor grado de interacción son más
predecibles" resultó ser exactamente la opuesta de lo medido ($r=-0.95$).
\end{abstract}

\section{Introducción}
La pregunta de si un algoritmo cuántico dado debe correr bajo QEM (mitigación,
sin qubits lógicos), QEC (corrección con overhead moderado) o FTQC completo
(overhead grande pero error arbitrariamente bajo) suele tratarse como una
propiedad fija del hardware disponible. En este trabajo mostramos que es, en
cambio, una función del circuito \emph{después} de aplicarle una técnica de
reducción --- y que la técnica elegida puede cambiar la decisión. Motivado por
esto, esta sesión de investigación evaluó de forma acumulativa nueve técnicas
de reducción/transformación de circuitos, siempre exigiendo (a) una
implementación real o una reproducción fiel del mecanismo publicado, (b)
verificación numérica cuando aplica, y (c) comparación cuantitativa contra
números del estado del arte cuando existen. Este artículo consolida esos
resultados.

\section{Técnicas evaluadas y estado del arte}
La Tabla~\ref{tab:sota} resume las nueve técnicas evaluadas, su fuente
publicada, y la herramienta de referencia usada para validarlas (propia o de
terceros).

\begin{table}[t]
\centering
\footnotesize
\caption{Técnicas evaluadas y su validación contra el estado del arte.}
\label{tab:sota}
\begin{tabular}{p{1.55cm}p{2.5cm}p{3.4cm}}
\toprule
\textbf{Técnica} & \textbf{Referencia} & \textbf{Validación} \\
\midrule
Cálculo-ZX & Duncan et al.\ 2020; PyZX & \texttt{full\_reduce()} real (PyZX) \\
Teletp.\ de puerta & Gottesman--Chuang 1999 & verificación numérica propia \\
Des-cómputo AND & Gidney 2018; Jones 2013 & verificación numérica propia (fid.=1.0) \\
QCP acotado & Chen et al., ICCS 2025 & reproducción propia (\texttt{qcp\_bounded\_tracker.py}) \\
BQCP & arXiv:2606.02018 (2026) & código real de terceros (GitHub) \\
Corte de circuitos & Qiskit \texttt{addon-cutting} & paquete real (\texttt{find\_cuts}) \\
RQAOA & Bravyi et al.\ 2019 & \texttt{OpenQAOA} real (Entropica Labs) \\
SwapStrategy & Qiskit \texttt{Commuting2qGate\allowbreak Router} & pase real de Qiskit \\
Azure QRE & Microsoft QRE & API real de estimación \\
\bottomrule
\end{tabular}
\end{table}

\section{Metodología: normalización al basis gates real}
Un error metodológico recurrente en este tipo de comparaciones es medir
T-count sobre una base de compuertas sintética (p.\ ej.\ \texttt{u3}), donde
una fase no-Clifford arbitraria queda oculta dentro de una compuerta genérica
y el T-count aparente cae a 0 sin que el circuito realmente lo permita.
Corregimos esto normalizando \emph{las tres técnicas} comparadas (sin
reducir, \texttt{transpile} nivel 3, cálculo-ZX) al mismo basis gates real de
un backend IBM Heron (\texttt{FakeTorino}: \texttt{cz, rz, sx, x}) antes de
perfilar cualquier métrica --- así \texttt{rz} con fase arbitraria es la
única fuente de no-Clifford, y es visible tanto para el perfilado como para
el cálculo posterior de error/tiempo reales de 2 qubits que usa
\texttt{decide\_regime()}.

La decisión de régimen compara, en la misma unidad de tiempo esperado:
\[
t_{\text{QEM}} = t_{\text{shot}} \cdot \frac{1}{F_{\text{esperada}}^2}
\qquad\text{vs.}\qquad
t_{\text{FTQC}} = \text{runtime}_{\text{Azure QRE}}
\]
sujeto a una restricción dura de factibilidad: si los qubits físicos que pide
FTQC exceden el presupuesto disponible, la decisión es QEM sin importar el
tiempo.

\section{Resultados}

\subsection{Caso límite: el cálculo-ZX puede hacer ganar a FTQC en tiempo}
La Figura~\ref{fig:regime} muestra el tiempo esperado (escala log) de QEM y
FTQC para las 4 familias de circuito, comparando la versión sin reducir
(\emph{baseline}) contra la versión reducida con cálculo-ZX. En la versión
baseline, QEM gana siempre por varios órdenes de magnitud. Tras cálculo-ZX,
FTQC pasa a ganar en \emph{tiempo} en las 4 familias (hasta 212$\times$ más
rápido en QAOA/farm) --- el único de los tres métodos de preparación probados
donde esto ocurre. Sin embargo, la restricción de presupuesto de qubits
físicos sigue bloqueando FTQC como opción viable en términos absolutos en
todos los casos: el resultado es un caso límite genuino (la ventaja de
tiempo existe) pero no una recomendación práctica de cambiar de régimen.

\begin{figure}[t]
\centering
\includegraphics[width=\linewidth]{figs/fig3_regime_case.pdf}
\caption{Tiempo esperado QEM vs.\ FTQC, antes y después de cálculo-ZX, para
4 familias de circuito (escala log).}
\label{fig:regime}
\end{figure}

\subsection{Validación cuantitativa contra el estado del arte}
La Figura~\ref{fig:lit} compara nuestras métricas medidas contra los rangos
publicados. La reducción de T-count por cálculo-ZX que medimos (16--53\%,
promedio 34.5\%) queda por debajo del rango publicado (51--75\%,
Kissinger--van~de~Wetering 2020), consistente con que \texttt{full\_reduce()}
por sí solo (sin post-procesamiento adicional de extracción de circuito
específico del \emph{paper}) es una cota inferior razonable. La reducción de
profundidad por teletransportación de puertas RZZ (17.5\% promedio) cae
dentro del rango publicado (10--25\%). El des-cómputo de puertas AND por
medición reproduce exactamente el 50\% de reducción de T-count reportado por
Gidney (2018) para el paso de des-cómputo, verificado numéricamente con
fidelidad exacta (20 estados aleatorios, ambas ramas de medición). El
incremento de conteo de compuertas de 2 qubits que observamos tras
cálculo-ZX (un efecto aparentemente contraintuitivo) está documentado de
forma independiente en la literatura (arXiv:2311.08881), por lo que no se
trata de un artefacto de nuestra implementación.

\begin{figure}[t]
\centering
\includegraphics[width=\linewidth]{figs/fig2_technique_vs_literature.pdf}
\caption{Comparación de nuestras métricas medidas contra rangos publicados.}
\label{fig:lit}
\end{figure}

\subsection{Las técnicas no siempre se combinan bien}
Se probó empíricamente si encadenar cálculo-ZX seguido de
teletransportación de puertas produce una reducción acumulada mayor que
aplicar cada técnica por separado. El resultado fue negativo: encadenarlas
infla el grado de interacción del grafo del circuito entre +35\% y +141\% en
3 de 4 familias, en vez de reducirlo. Esto valida (post-hoc, no por diseño
original) la arquitectura de \texttt{pipeline\_v2.py}, que ya mantenía ramas
separadas por régimen en vez de una cadena de técnicas.

\subsection{Propagación de constantes cuántica (QCP/BQCP): resultado negativo}
Implementamos una versión acotada de \emph{Quantum Constant Propagation}
(Chen et al., ICCS 2025) con un tamaño máximo de grupo rastreado
($n_{\text{pcm}}$) y un atajo de "control clásico" para evitar inflar grupos
innecesariamente. Sobre nuestras 4 familias de circuito (con entradas puestas
en superposición genuina vía H+S, evitando el caso degenerado de entrada
clásica), la fracción de qubits con estado conocido permanece en 0\% para
$n_{\text{pcm}} \le 8$--16, y solo salta a 100\% cuando el grupo rastreado
cubre el circuito completo --- es decir, sin una utilidad práctica intermedia
en estos circuitos. Validamos este hallazgo negativo con el código real de
BQCP (\emph{Branch-Aware QCP}, arXiv:2606.02018), que en nuestros 4 circuitos
estáticos obtiene 0.0\%, 0.3\%, 0.0\% y 2.2\% de compuertas eliminadas,
frente a 3.6--4.3\% en su propio generador de circuitos dinámicos (para el
que fue diseñado). La Figura~\ref{fig:bqcp} muestra ambos conjuntos.

\begin{figure}[t]
\centering
\includegraphics[width=\linewidth]{figs/fig4_bqcp.pdf}
\caption{Reducción de compuertas por BQCP: circuitos estáticos propios (rojo)
vs.\ circuitos dinámicos de BQCP (azul).}
\label{fig:bqcp}
\end{figure}

\subsection{Corte de circuitos: viable solo si el grafo es separable}
Usando el paquete real \texttt{qiskit-addon-cutting}, encontramos que el
\emph{overhead} de muestreo depende críticamente de si el grafo de
interacción del circuito admite un corte de tamaño pequeño. Sumador8 (grafo
en cadena, altamente separable) necesita solo 2 cortes con un
\emph{overhead} de $256\times$ --- tratable. QAOA/farm (grafo denso de 17
qubits, 39 aristas) requiere 18--44 cortes según la estrategia, con
\emph{overhead} de $10^{17}$--$10^{43}\times$ --- astronómicamente
intratable. La conclusión es que la \emph{separabilidad estructural} del
grafo de interacción, no solo el número de qubits, es el predictor correcto
de si el corte de circuitos es viable.

\subsection{RQAOA: la intuición sobre "qubits de control" era incorrecta}
Se propuso la hipótesis de que las aristas con mayor grado de interacción
(qubits más "centrales") tendrían correlaciones $\langle Z_iZ_j\rangle$ más
confiables/deterministas, permitiendo tratarlas como controles fijos. Usando
\texttt{OpenQAOA} (RQAOA real, Bravyi et al.\ 2019) medimos esto
directamente: la correlación entre la suma de grados de una arista y la
magnitud de $|\langle Z_iZ_j\rangle|$ resultó fuertemente
\emph{negativa} ($r=-0.95$, Figura~\ref{fig:rqaoa}) --- exactamente lo
opuesto de la hipótesis. Las aristas de \emph{bajo} grado son las más
confiables de fijar primero, no las de alto grado. Reportamos esto como
resultado empírico honesto, no como confirmación de la idea original.

\begin{figure}[t]
\centering
\includegraphics[width=\linewidth]{figs/fig5_rqaoa_correlation.pdf}
\caption{Correlación $|\langle Z_iZ_j\rangle|$ vs.\ suma de grados de la
arista en QAOA/farm ($r=-0.95$).}
\label{fig:rqaoa}
\end{figure}

\subsection{SwapStrategy vs.\ SabreSwap: un cruce dependiente de densidad}
Se comparó el enrutamiento estructurado de Qiskit
(\texttt{SwapStrategy}/\texttt{Commuting2qGateRouter}, para circuitos con
compuertas de 2 qubits mutuamente conmutables como QAOA) contra el
enrutamiento heurístico genérico \texttt{SabreSwap}, sobre grafos aleatorios
de 17 qubits con densidad creciente. El costo de \texttt{SwapStrategy} se
mantiene esencialmente constante ($\sim$120 SWAPs, profundidad 54--63)
independientemente de la densidad, mientras que \texttt{SabreSwap} escala
con la densidad (Figura~\ref{fig:density}). El cruce ocurre antes en
profundidad (densidad~$\approx$0.30) que en conteo de SWAPs
(densidad~$\approx$0.45--0.60): para grafos suficientemente densos,
\texttt{SwapStrategy} gana en ambas métricas; para grafos dispersos,
\texttt{SabreSwap} es mejor en las dos.

Nota de honestidad metodológica: se observa una caída anómala en el conteo
de SWAPs de \texttt{SabreSwap} entre densidad 0.80 (274 SWAPs) y densidad
1.00/grafo completo (132 SWAPs) --- lo contrario de lo esperado si la
tendencia fuera monótona. No se ha investigado la causa raíz de esta caída;
una hipótesis plausible es que un grafo completo, al ser totalmente
simétrico, elimina la variabilidad estructural que hace difícil a la
heurística de Sabre encontrar una solución cercana al óptimo, pero esto no
se ha verificado y se reporta como observación abierta, no como conclusión.

\begin{figure}[t]
\centering
\includegraphics[width=\linewidth]{figs/fig1_density_sweep.pdf}
\caption{SWAPs (a) y profundidad (b) vs.\ densidad del grafo, para
\texttt{SwapStrategy} (constante) y \texttt{SabreSwap} (creciente, con una
anomalía en densidad=1.0).}
\label{fig:density}
\end{figure}

\section{Discusión}
Un patrón transversal emerge de las nueve técnicas: la utilidad de cada una
depende de una propiedad \emph{estructural} medible del circuito o de su
grafo de interacción --- no de su tamaño en qubits o compuertas por sí solo.
El cálculo-ZX ayuda cuando el T-count relativo es alto; el des-cómputo por
medición ayuda cuando existen ancillas de vida corta con un patrón AND
explícito; el corte de circuitos ayuda cuando el grafo admite cortes
pequeños (separabilidad); RQAOA se beneficia de fijar primero las aristas de
\emph{bajo} grado, no de alto grado; y \texttt{SwapStrategy} gana quando la
densidad de interacción es alta. Esto motiva directamente el diseño de
\texttt{technique\_advisor.py}: un evaluador automático que mide estas
propiedades estructurales (grado promedio/máximo, fracción de T-gates,
candidatos a ancilla, fracción conocida de QCP) y emite recomendaciones
SI/NO calibradas empíricamente en esta sesión, en vez de aplicar todas las
técnicas indiscriminadamente.

\section{Conclusión}
Este trabajo no reivindica superioridad de ninguna técnica sobre el estado
del arte publicado; de hecho, en tres métricas comparables nuestras cifras
quedan dentro o por debajo de los rangos publicados. La contribución es una
capa de \emph{decisión} que: (1) corrige un error metodológico
recurrente (T-count oculto por normalización a basis gates sintéticos); (2)
cuantifica cuándo la elección de técnica de reducción cambia la decisión
QEM/FTQC; (3) identifica qué propiedad estructural del circuito predice la
utilidad de cada técnica, validado empíricamente y no solo por intuición
--- incluyendo dos casos donde la intuición inicial resultó incorrecta
(encadenamiento de técnicas, y grado de qubit en RQAOA); y (4) deja
documentada, sin ocultar, una anomalía sin explicar (caída de SWAPs de Sabre
en el grafo completo) para trabajo futuro.

\section*{Reproducibilidad}
Todo el código, datos y figuras de este artículo provienen de ejecuciones
reales sobre PyZX, Qiskit, \texttt{qiskit-ibm-runtime} (\texttt{FakeTorino}),
\texttt{qiskit-addon-cutting}, OpenQAOA, y el código público de BQCP
(\url{https://github.com/1nnocenzo/bqcp}). Ningún número reportado es
estimado a mano.

\begin{thebibliography}{9}
\footnotesize
\bibitem{kissinger2020}
R.\ Kissinger, J.\ van~de~Wetering, ``Reducing T-count with the ZX-calculus,''
\emph{Phys.\ Rev.\ A}, 102, 022406, 2020.

\bibitem{gidney2018}
C.\ Gidney, ``Halving the cost of quantum addition,''
\emph{Quantum} 2, 74, 2018.

\bibitem{jones2013}
C.\ Jones, ``Low-overhead constructions for the fault-tolerant Toffoli gate,''
\emph{Phys.\ Rev.\ A} 87, 022328, 2013.

\bibitem{chen2025}
Chen, Fulginiti, Mendl, ``Optimization Framework for Reducing Mid-circuit
Measurements and Resets,'' ICCS 2025.

\bibitem{bqcp2026}
``Branch-Aware Quantum Constant Propagation,'' arXiv:2606.02018, 2026.
\url{https://github.com/1nnocenzo/bqcp}

\bibitem{bravyi2019}
S.\ Bravyi, A.\ Kliesch, R.\ Koenig, E.\ Tang, ``Obstacles to Variational
Quantum Optimization from Symmetry Protection,'' \emph{Phys.\ Rev.\ Lett.}
125, 260505, 2020 (RQAOA).

\bibitem{gottesman1999}
D.\ Gottesman, I.\ Chuang, ``Demonstrating the viability of universal
quantum computation using teleportation and single-qubit operations,''
\emph{Nature} 402, 390--393, 1999.

\bibitem{qoblib}
ZIB, ``QOBLIB: Quantum Optimization Benchmark Library,''
\url{https://github.com/ZIB-AOPT/QOBLIB}.

\bibitem{azureqre}
Microsoft, ``Azure Quantum Resource Estimator,'' documentación técnica,
2024.
\end{thebibliography}

\end{document}
